\documentclass[10pt,a4paper]{amsart}
\usepackage[margin=19mm]{geometry}
\usepackage{amsmath,amssymb,mathtools}
\usepackage[scr=rsfs,cal=euler]{mathalfa}
\usepackage{xfrac}
\usepackage[unicode,hidelinks,bookmarks=false]{hyperref}
\newcommand{\ie}{i.e.\ }
\newcommand{\cf}{cf.\ }
\newcommand{\resp}{respectively\ }
\newcommand{\Subsection}{\subsection}
\providecommand{\nobreakdash}{\nobreak\hbox{-}}
\setlength{\emergencystretch}{2em}
\multlinegap0pt
\usepackage{bm}
\usepackage{bbm}
\usepackage{dsfont}
\usepackage{xspace}
\usepackage{cancel}
\usepackage{mathtools}
\usepackage[shortlabels]{enumitem}
\usepackage{amsthm}
\makeatletter
\@ifundefined{theorem}{\newtheorem{theorem}{Theorem}}{}
\makeatother
\title[An infinite family of internal congruences modulo powers of ]{An infinite family of internal congruences modulo powers of 2 for partitions into odd parts with designated summands}
\author{Shane Chern, James A. Sellers}
\date{Source: arXiv:2308.04348v1 (CC BY 4.0)}
\begin{document}
\maketitle

\noindent\emph{Source and licence.} An infinite family of internal congruences modulo powers of 2 for partitions into odd parts with designated summands. Version arXiv:2308.04348v1 (CC BY 4.0), distributed under the Creative Commons Attribution 4.0 International licence, as recorded in the arXiv licence metadata for this exact version and shown on the abstract page https://arxiv.org/abs/2308.04348v1. Full rights notice: licenses/arxiv-2308.04348-CC-BY-4.0.txt.\par\smallskip

\noindent\textit{Editorial scope.} Counted unit: the Theorem whose source label is \texttt{th:F-nonvanishing} in the article (source0.tex lines 854--866, source label \texttt{th:F-nonvanishing}), delivered together with the article's own complete original proof of it (source0.tex lines 973--1096, one \texttt{proof} environment of 124 lines). No mathematics is written, repaired or re-derived: statement and proof are the article's own text line for line. The proof is detached from the statement in the source: the article states the result at lines 854--866 and prints its proof at lines 973--1096, and the proof environment's own header names the statement, so the association is the article's own. Required context retained uncounted: th:Z-nonvanishing (lines 467--482); th:Z-2-adic (lines 702--718); eq:F-xi-rec (lines 835--837); eq:tau-even-mod-4 (lines 870--873); eq:tau-odd-mod-4 (lines 870--873); eq:F-nu-m (lines 958--960). Every definition, display and label of those lines is retained unchanged. Omitted from this package and not redistributed: 1--466, 483--701, 719--834, 838--853, 867--869, 874--957, 961--972, 1097--1238. Nothing inside a retained excerpt is omitted or altered: every retained body is delivered exactly as the article's lines are written, so no omission receipt of any kind accompanies this package. Citations: the retained excerpts carry no citation command at all, so no bibliography entry of the article is transcribed and no reference is redistributed. Reference adaptation: none was needed and none was made; every reference inside the delivered unit resolves inside it, so no unresolved reference or citation is produced. Printed slips kept literal: no printed word, symbol, hypothesis or number of the counted statement or of its proof was altered. Layout and class: the article's own class is \texttt{amsart} with packages including tikz, xcolor, latexsym, amsmath, amssymb, epic, amsthm, color, geometry, etoolbox, enumitem, cite, hyperref; the wrapper is the lane's pinned \texttt{amsart} editorial wrapper at 10pt on A4, which restates the theorem-like environments the retained text needs and adds the article's own macro definitions that the retained text needs (none), with extra wrapper packages (bm, bbm, dsfont, xspace, cancel, mathtools, enumitem). The delivered numbering is the unit's own. Assets: no retained span contains a figure environment, a graphics-inclusion command, a PDF-page inclusion, a table or a TikZ picture; no image file is copied into this package, the package declares no asset dependency and no image file is redistributed with it. Licence: version arXiv:2308.04348v1 of this article is distributed under the Creative Commons Attribution 4.0 International licence, as recorded in the arXiv licence metadata for this exact version and shown on the abstract page https://arxiv.org/abs/2308.04348v1. A full rights notice is delivered at licenses/arxiv-2308.04348-CC-BY-4.0.txt. Characters: every retained excerpt is pure ASCII, so the delivered engine, XeLaTeX with its default Latin Modern text font, reports no missing character.
\begin{theorem}\label{th:Z-nonvanishing}
	For any $i,j\ge 0$, define
	\begin{align*}
		d_{i,j}:=\begin{cases}
			5I+J, & \text{if $(i,j)=(2I,2J)$},\\
			5I+J+1, & \text{if $(i,j)=(2I,2J+1)$},\\
			5I+J+3, & \text{if $(i,j)=(2I+1,2J)$},\\
			5I+J+3, & \text{if $(i,j)=(2I+1,2J+1)$}.
		\end{cases}
	\end{align*}
	Then for any $m$ with $0\le m< d_{i,j}$, we have
	\begin{align*}
		Z_{i,j}(m) = 0.
	\end{align*}
	Furthermore, the coefficient $Z_{i,j}(d_{i,j})$ is an odd integer.
\end{theorem}

\begin{theorem}\label{th:Z-2-adic}
	For any $k\ge 2$ and $j\ge 0$, we have
	\begin{align*} %\label{eq:Z-2-adic-0}
		\nu\big(Z_{2^k,j}(d_{2^k,j})\big) = 0,
	\end{align*}
	and
	\begin{align*} %\label{eq:Z-2-adic-1}
		\nu\big(Z_{2^k,j}(d_{2^k,j}+1)\big) \begin{cases}
			= 1, & \text{if $j\equiv 2 \bmod{4}$},\\
			\ge 2, & \text{if $j\not\equiv 2 \bmod{4}$}.
		\end{cases}
	\end{align*}
	Furthermore, for $M\ge 2$,
	\begin{align*} %\label{eq:Z-2-adic-M}
		\nu\big(Z_{2^k,j}(d_{2^k,j}+M)\big) \ge M+1.
	\end{align*}
\end{theorem}

	\begin{align}\label{eq:F-xi-rec}
		\Phi_k = \sum_{\ell} F_{k-1}(\ell) \cdot \zeta_{2^{k-1},\ell},
	\end{align}

	\begin{align}
		\tau_{2K-1} &\equiv 2 \pmod{4},\label{eq:tau-odd-mod-4}\\
		\tau_{2K} &\equiv 3 \pmod{4},\label{eq:tau-even-mod-4}
	\end{align}

	\begin{align}\label{eq:F-nu-m}
		\nu\big(F_{2K+1}(\tau_{2K+1}+M)\big) &\ge 2K+M+2.
	\end{align}

\begin{theorem}\label{th:F-nonvanishing}
	For any $k\ge 3$, define
	\begin{align*}
		\tau_{k}:=\begin{cases}
			7\cdot 2^{2K-3}-\tfrac{2}{3}\big(4^{K-2}-1\big), & \text{if $k=2K-1$},\\[6pt]
			7\cdot 2^{2K-2}-\tfrac{1}{3}\big(4^{K-1}-1\big), & \text{if $k=2K$}.
		\end{cases}
	\end{align*}
	Then for any $m$ with $0\le m< \tau_k$, we have
	\begin{align*}
		F_k(m) = 0.
	\end{align*}
\end{theorem}

\begin{proof}[Proof of \eqref{eq:F-nu-m}]
	To produce $\Phi_{2K+1}$ from $\Phi_{2K-1}$, we need to employ the recurrence \eqref{eq:F-xi-rec} twice. Writing
	\begin{align*}
		\Phi_{2K-1} = \sum_{M''\ge 0} F_{2K-1}(\tau'' + M'')\xi^{\tau''+M''},
	\end{align*}
	we deduce from \eqref{eq:F-xi-rec} that
	\begin{align*}
		\Phi_{2K} = \sum_{M''\ge 0} F_{2K-1}(\tau'' + M'')\zeta_{2^{2K-1},\tau'' + M''}.
	\end{align*}
	Recall that $\tau''$ is even according to \eqref{eq:tau-odd-mod-4}. It then follows from Theorems \ref{th:Z-nonvanishing} and \ref{th:F-nonvanishing} that the minimal degree of the $\xi$-powers in $\zeta_{2^{2K-1},\tau'' + M''}$ is $\tau'+\lfloor\tfrac{M''+1}{2}\rfloor$, i.e.,
	\begin{align}\label{eq:d-tau-2}
		d_{2^{2K-1},\tau'' + M''} = \tau'+\lfloor\tfrac{M''+1}{2}\rfloor.
	\end{align}
	Hence,
	\begin{align*}
		\Phi_{2K} = \sum_{\substack{M''\ge 0\\M'\ge \lfloor\frac{M''+1}{2}\rfloor}} F_{2K-1}(\tau'' + M'')Z''(\tau'' + M'',\tau'+M')\xi^{\tau'+M'}.
	\end{align*}
	Applying the recurrence \eqref{eq:F-xi-rec} once again gives us that
	\begin{align*}
		\Phi_{2K+1} = \sum_{\substack{M''\ge 0\\M'\ge \lfloor\frac{M''+1}{2}\rfloor}} F_{2K-1}(\tau'' + M'')Z''(\tau'' + M'',\tau'+M')\zeta_{2^{2K},\tau' + M'}.
	\end{align*}
	Noting that $\tau'$ is odd according to \eqref{eq:tau-even-mod-4}, Theorems \ref{th:Z-nonvanishing} and \ref{th:F-nonvanishing} then imply that the minimal degree of the $\xi$-powers in $\zeta_{2^{2K},\tau' + M'}$ is $\tau+\lfloor\tfrac{M'}{2}\rfloor$, i.e.,
	\begin{align}\label{eq:d-tau-1}
		d_{2^{2K},\tau' + M'} = \tau'+\lfloor\tfrac{M'}{2}\rfloor.
	\end{align}
	Hence,
	\begin{align*}
		\Phi_{2K+1} = \sum_{\substack{M''\ge 0\\M'\ge \lfloor\frac{M''+1}{2}\rfloor\\M\ge \lfloor\frac{M'}{2}\rfloor}} C(M,M'',M'') \xi^{\tau+M},
	\end{align*}
	where
	\begin{align*}
		C(M,M'',M''):=F_{2K-1}(\tau'' + M'')Z''(\tau'' + M'',\tau'+M')Z'(\tau' + M',\tau+M).
	\end{align*}
	
	Note that to complete the inductive step, it is sufficient to show the inequality
	\begin{align}\label{eq:nu-C}
		\nu\big(C(M,M'',M'')\big) \ge 2K+M+2
	\end{align}
	for every tuple $(M,M'',M'')$ such that $M''\ge 0$, $M'\ge \lfloor\tfrac{M''+1}{2}\rfloor$ and $M\ge \lfloor\tfrac{M'}{2}\rfloor$.
	
	We first consider the following two generic cases where \textbf{(i).}~$M''\ge 2$ and \textbf{(ii).}~$M''=0$ or $1$ and $M'\ge 3$. By our inductive assumption,
	\begin{align*}
		\nu\big(F_{2K-1}(\tau'' + M'')\big) \ge 2K + M''.
	\end{align*}
	Also, by a weaker form of Theorem \ref{th:Z-2-adic} with recourse to \eqref{eq:d-tau-2} and \eqref{eq:d-tau-1}, respectively, we have
	\begin{align*}
		\nu\big(Z''(\tau'' + M'',\tau'+M')\big) &\ge M'-\lfloor\tfrac{M''+1}{2}\rfloor,\\
		\nu\big(Z'(\tau' + M',\tau+M)\big) &\ge M-\lfloor\tfrac{M'}{2}\rfloor.
	\end{align*}
	It follows that
	\begin{align*}
		\nu\big(C(M,M'',M'')\big) \ge 2K+M+\big(M''-\lfloor\tfrac{M''+1}{2}\rfloor\big) + \big(M'-\lfloor\tfrac{M'}{2}\rfloor\big).
	\end{align*}
	Now \textbf{(i).}~if $M''\ge 2$, then $M'\ge 1$ so that $M''-\lfloor\tfrac{M''+1}{2}\rfloor\ge 1$ and that $M'-\lfloor\tfrac{M'}{2}\rfloor\ge 1$; \textbf{(ii).}~if $M''=0$ or $1$ and $M'\ge 3$, then $M''-\lfloor\tfrac{M''+1}{2}\rfloor= 0$ but $M'-\lfloor\tfrac{M'}{2}\rfloor\ge 2$. Thus, in both cases the inequality \eqref{eq:nu-C} holds.
	
	Now there are five sporadic cases to be investigated, and we work on them separately. Note that for \eqref{eq:F-nu-m}, we have the condition that $M\ge 1$.
	\begin{enumerate}[label={\textbf{(\alph*).}},leftmargin=*,labelsep=0cm,align=left]
		\item $(M',M'')=(0,0)$: In this case, it is known by the inductive assumption,
		\begin{align*}
			\nu\big(F_{2K-1}(\tau'' + 0)\big) \ge 2K + 1,
		\end{align*}
		that by Theorem \ref{th:Z-2-adic},
		\begin{align*}
			\nu\big(Z''(\tau'' + 0,\tau'+0)\big) = 0,
		\end{align*}
		and that also by Theorem \ref{th:Z-2-adic} (noting that $\tau'+0\equiv 3\pmod{4}$),
		\begin{align*}
			\nu\big(Z'(\tau' + 0,\tau+M)\big) \ge M+1.
		\end{align*}
	
		\item $(M',M'')=(1,0)$: In this case, it is known by the inductive assumption,
		\begin{align*}
			\nu\big(F_{2K-1}(\tau'' + 0)\big) \ge 2K + 1,
		\end{align*}
		that by Theorem \ref{th:Z-2-adic} (noting that $\tau''+0\equiv 2 \pmod{4}$),
		\begin{align*}
			\nu\big(Z''(\tau'' + 0,\tau'+1)\big) = 1,
		\end{align*}
		and that also by Theorem \ref{th:Z-2-adic} (noting that $\tau'+1\equiv 0\pmod{4}$),
		\begin{align*}
			\nu\big(Z'(\tau' + 1,\tau+M)\big) \ge M+1.
		\end{align*}
	
		\item $(M',M'')=(2,0)$: In this case, it is known by the inductive assumption,
		\begin{align*}
			\nu\big(F_{2K-1}(\tau'' + 0)\big) \ge 2K + 1,
		\end{align*}
		that by Theorem \ref{th:Z-2-adic},
		\begin{align*}
			\nu\big(Z''(\tau'' + 0,\tau'+2)\big) \ge 3,
		\end{align*}
		and that by a weaker form of Theorem \ref{th:Z-2-adic},
		\begin{align*}
			\nu\big(Z'(\tau' + 2,\tau+M)\big) \ge M-1.
		\end{align*}
	
		\item $(M',M'')=(1,1)$: In this case, it is known by the inductive assumption,
		\begin{align*}
			\nu\big(F_{2K-1}(\tau'' + 1)\big) \ge 2K + 1,
		\end{align*}
		that by Theorem \ref{th:Z-2-adic},
		\begin{align*}
			\nu\big(Z''(\tau'' + 1,\tau'+1)\big) = 0,
		\end{align*}
		and that also by Theorem \ref{th:Z-2-adic} (noting that $\tau'+1\equiv 0\pmod{4}$),
		\begin{align*}
			\nu\big(Z'(\tau' + 1,\tau+M)\big) \ge M+1.
		\end{align*}
	
		\item $(M',M'')=(2,1)$: In this case, it is known by the inductive assumption,
		\begin{align*}
			\nu\big(F_{2K-1}(\tau'' + 1)\big) \ge 2K + 1,
		\end{align*}
		that by Theorem \ref{th:Z-2-adic} (noting that $\tau''+1\equiv 3 \pmod{4}$),
		\begin{align*}
			\nu\big(Z''(\tau'' + 1,\tau'+2)\big) \ge 2,
		\end{align*}
		and that by a weaker form of Theorem \ref{th:Z-2-adic},
		\begin{align*}
			\nu\big(Z'(\tau' + 2,\tau+M)\big) \ge M-1.
		\end{align*}
	\end{enumerate}
	We find that in each of these cases, the inequality \eqref{eq:nu-C} is also valid, thereby closing our proof of \eqref{eq:F-nu-m}.
\end{proof}

\end{document}
